% Plan de mise à niveau d'EP Manager (mod_ep + mod_epsynthesis) au niveau de StageCompagnon.
% Compilation : latexmk -pdf plan-mise-a-niveau.tex
\documentclass[11pt,a4paper]{article}
\usepackage[french]{babel}
\usepackage[T1]{fontenc}
\usepackage[utf8]{inputenc}
\usepackage{lmodern,microtype,booktabs,tabularx,longtable,xcolor,enumitem,hyperref}
\usepackage[margin=2.2cm]{geometry}
\hypersetup{colorlinks=true, linkcolor=blue!50!black, urlcolor=blue!50!black}
\newcommand{\code}[1]{\texttt{\detokenize{#1}}}
\newcommand{\lot}[1]{\textbf{Lot~#1}}
\setlist[itemize]{nosep}
\setlist[enumerate]{nosep}

\title{EP Manager : plan de mise à niveau\\
  \large Enseignement personnalisé sous Moodle, au niveau de StageCompagnon}
\author{Sébastien Lefebvre}
\date{6 novembre 2026 -- version 1 du plan, soumise à validation}

\begin{document}
\maketitle

\section{Objet}

Ce document fixe ce que la mise à niveau d'EP Manager doit apporter, dans quel ordre et
avec quelles vérifications, avant toute ligne de code. Il part de l'état actuel des deux
dépôts (\code{EPmanager}, version \code{0.1.0} de \code{mod_ep} et
\code{mod_epsynthesis}~; \code{StageCompagnon}, modules \code{mod_stage} et
\code{mod_stagesynthesis}), de la demande du 6~novembre~2026 et des références publiques
sur le supplément au diplôme et la description des unités d'enseignement.

Les attendus de la demande sont repris un par un en section~\ref{sec:cible}. La
section~\ref{sec:lots} les découpe en lots livrables séparément, chacun avec ses tests. La
section~\ref{sec:questions} liste les hypothèses prises pour avancer sans attendre et les
quelques points qui méritent une décision.

Principe retenu partout : tout ce qui relève de l'administration se fait depuis le panneau
\textbf{Administration} de l'activité \code{mod_ep}, par la DEVE (capacité
\code{mod/ep:manage}). Rien ne demande un rôle d'administrateur du site ni un passage par
\emph{Administration du site}, en dehors de l'installation du plugin et de la création
des rôles standard déjà décrite dans StageCompagnon.

\section{État des lieux}

\subsection{EP Manager aujourd'hui}

\code{mod_ep} (environ 7\,000 lignes de PHP, dont 2\,100 dans \code{locallib.php}) gère~:
six types d'EP fixes (académique interne, stage, engagement étudiant, expérience
professionnelle, sport, académique externe), un catalogue d'EP internes avec places,
plage d'années et responsables, l'inscription de l'étudiant au catalogue, la déclaration
hors catalogue avec justificatifs, la validation par le responsable de l'EP ou par
l'enseignant référent (lu dans \code{mod_stage}) ou la DEVE, l'attribution automatique
depuis les stages complémentaires validés, les plafonds par type (cursus, année), les
minimums par année et de cursus, un tableau de pilotage, deux exports CSV. Quatre
fichiers de tests PHPUnit (842 lignes) couvrent le décompte, le catalogue et la
synchronisation des stages. \code{mod_epsynthesis} (1\,200 lignes) agrège la validation et le
pilotage pour un enseignant, toutes promotions confondues.

\subsection{Ce que StageCompagnon a et qu'EP Manager n'a pas}

\begin{longtable}{p{4.2cm}p{5.3cm}p{5.3cm}}
\toprule
\textbf{Sujet} & \textbf{StageCompagnon} & \textbf{EP Manager} \\
\midrule
\endhead
Sauvegarde Moodle & \code{backup/moodle2/} complet, données utilisateur et fichiers, tests de
  restauration & \code{FEATURE_BACKUP_MOODLE2} annoncé mais \textbf{aucune implémentation}~:
  une sauvegarde de cours perd tout \\
Journaux & Un événement Moodle par transition du circuit & Aucun événement \\
Gabarits & Mustache et classes \code{output} pour les listes & Tout en \code{html_writer} \\
Imports & CSV générique, StageVet, historique, mobilités, restauration d'export global,
  import depuis une autre activité, prévisualisation, verrou, transaction & Aucun import \\
Exports & Excel multi-feuilles (bilan, détail, réponses), PDF de bilan, services web & Deux CSV \\
Concurrence & Verrou par activité, écritures conditionnées à l'état lu, audit documenté &
  Aucune protection (voir §\ref{sec:perf}) \\
Charge & Pagination SQL, cache de bilan de promotion, chargements groupés, banc de charge &
  Pagination en PHP après chargement complet, N requêtes par étudiant \\
Tests & 45 fichiers PHPUnit, Behat, CI sur trois combinaisons, couverture & 4 fichiers, pas de
  CI, pas de Behat \\
Documentation & README, INSTALL détaillé, CHANGELOG, trois guides LaTeX avec captures,
  rapports d'audit & README et INSTALL courts \\
Moodle minimal & 4.3, PHP 8.0, justifié & 4.0 annoncé, non vérifié \\
\bottomrule
\end{longtable}

\subsection{Défauts repérés à la lecture}\label{sec:defauts}

Ces points sont corrigés dans les lots, mais sont listés ici parce qu'ils existent déjà en
production potentielle~:
\begin{enumerate}
  \item \textbf{Places du catalogue sous concurrence.} Le contrôle « reste-t-il une place »
    puis l'insertion ne sont pas atomiques~: deux étudiants qui confirment ensemble sur la
    dernière place sont tous deux inscrits.
  \item \textbf{Synchronisation des stages en double.} \code{ep_sync_stage_credits_if_due()}
    est appelée à l'affichage de chaque page~; deux pages ouvertes en même temps après
    cinq minutes lancent deux synchronisations qui créent chacune les crédits manquants (pas
    d'index unique sur \code{epid, source, sourceref}).
  \item \textbf{Décisions écrasées.} Validation, refus et annulation réécrivent la ligne sans
    condition sur l'état lu~: un responsable et la DEVE qui statuent ensemble se recouvrent,
    comme StageCompagnon l'a constaté sur les conventions avant correction.
  \item \textbf{Sauvegarde absente} (ci-dessus)~: une duplication de cours ou une
    restauration produit une activité sans types, sans catalogue, sans crédits.
  \item \textbf{Pilotage non paginé en base}~: \code{ep_get_pilotage_overview()} calcule le
    bilan de chaque étudiant inscrit (3 à 4 requêtes par étudiant) avant de filtrer, trier et
    découper en PHP~; l'export CSV refait de même. Les places restantes et les responsables
    du catalogue sont lus en une requête par EP.
  \item \textbf{Confirmation JavaScript non échappée}~: \code{onclick="return
    confirm('...')"} reçoit une chaîne de langue brute~; une apostrophe dans la traduction
    casse la page. À remplacer par les confirmations Moodle (\code{confirm_action} ou page de
    confirmation).
  \item \textbf{Type d'EP d'un crédit à l'édition}~: l'étudiant qui modifie une déclaration
    peut la basculer sur n'importe quel type déclarable, y compris après que la DEVE a
    rédigé une consigne spécifique~; acceptable, mais à journaliser.
\end{enumerate}

\section{Références externes}

\subsection{Supplément au diplôme}

Le supplément au diplôme est l'annexe descriptive prévue par le décret du 8~avril~2002
(code de l'éducation, art.~D.~123-13), délivrée gratuitement avec le diplôme, selon le
modèle commun Commission européenne / Conseil de l'Europe / UNESCO-CEPES repris par
Europass. Il ne remplace pas le diplôme et ne vaut pas reconnaissance. Sa structure est
fixe, en huit rubriques~:
\begin{enumerate}
  \item \textbf{Informations sur le titulaire}~: 1.1 nom(s) de famille, 1.2 prénom(s),
    1.3 date de naissance, 1.4 numéro ou code d'identification de l'étudiant.
  \item \textbf{Informations sur le diplôme}~: 2.1 intitulé du diplôme et titre conféré,
    2.2 principal(aux) domaine(s) d'étude, 2.3 nom et statut de l'établissement ayant
    délivré le diplôme, 2.4 nom et statut de l'établissement dispensant les cours,
    2.5 langue(s) d'enseignement et d'examen.
  \item \textbf{Niveau et durée}~: 3.1 niveau du diplôme, 3.2 durée officielle du
    programme (crédits, années), 3.3 conditions d'accès.
  \item \textbf{Programme suivi et résultats}~: 4.1 organisation des études, 4.2 acquis
    d'apprentissage du programme, \textbf{4.3 détail du programme, crédits obtenus et
    notes}, 4.4 système de notation et répartition des notes, 4.5 classification générale.
  \item \textbf{Fonction du diplôme}~: 5.1 accès à un niveau d'études supérieur,
    5.2 statut professionnel (profession réglementée).
  \item \textbf{Informations complémentaires}~: \textbf{6.1 informations complémentaires},
    6.2 autres sources d'information.
  \item \textbf{Certification du supplément}~: 7.1 date, 7.2 signature, 7.3 qualité du
    signataire, 7.4 cachet ou sceau officiel.
  \item \textbf{Système national d'enseignement supérieur} (texte national fourni par le
    ministère).
\end{enumerate}
Deux rubriques reçoivent ce que ce module produit~: la \textbf{4.3} accueille les EP
validés (intitulé, année, ECTS, note ou « validé »), la \textbf{6.1} les éléments que
l'étudiant demande à y voir figurer (engagements, distinctions, mobilités, certifications)
une fois validés. Les autres rubriques sont du texte d'établissement, saisi une fois par la
DEVE et transférable d'une activité à l'autre, sauf la rubrique~1 (profil de l'étudiant)
et la rubrique~7 (signataire et date, renseignés à l'émission).

Sources~: Europass, \emph{Supplément au diplôme}
(\url{https://europass.europa.eu/fr/what-europass-1/education-et-formation/supplement-au-diplome})~;
Agence Erasmus+ France, \emph{Supplément au diplôme : mode d'emploi} et trame officielle
(\url{https://agence.erasmusplus.fr/programme-erasmus/outils/europass/supplement-au-diplome/},
\url{https://agence.erasmusplus.fr/publications/guide-pratique-pour-la-mise-en-place-du-supplement-au-diplome/})~;
annexe~III de la décision Europass sur EUR-Lex
(\url{https://eur-lex.europa.eu/resource.html?uri=cellar:2ddb7ba0-339e-467d-ba4a-008837e1f1cb.0007.02/DOC_190&format=PDF})~;
ENIC-NARIC France, \emph{Le supplément au diplôme}
(\url{https://www.enic-naric.net/fileusers/8601_LE_SUPPLEMENT_AU_DIPLOME.pdf})~;
Euroguidance France, \emph{Les outils de la reconnaissance}
(\url{https://www.euroguidance-france.org/wp-content/uploads/2017/11/Les-outils-de-la-reconnaissance.pdf}).

\subsection{Fiche descriptive d'une unité d'enseignement (syllabus)}

Le guide d'utilisation ECTS (2015) prévoit, pour chaque unité du catalogue de cours, une
description d'une page au plus avec~: code, intitulé, type d'unité, niveau, année
d'étude, semestre, crédits ECTS, enseignant(s), acquis d'apprentissage, mode
d'enseignement, prérequis et corequis, contenu, lectures et ressources recommandées,
activités pédagogiques et méthodes, langue d'enseignement, modalités et critères
d'évaluation, charge de travail (un crédit valant environ 25 à 30~heures de travail de
l'étudiant). Les fiches francophones (ULB, UCLouvain, Fédération Wallonie-Bruxelles)
reprennent ces rubriques sous les intitulés « acquis d'apprentissage », « prérequis »,
« contenu », « méthodes d'enseignement », « modalités d'évaluation », « charge de travail
», « bibliographie »~; c'est ce vocabulaire qui sera employé dans l'interface.

Sources~: ECTS Users' Guide 2015
(\url{https://ehea.info/media.ehea.info/file/ECTS_Guide/77/4/ects-guide_en_595774.pdf})~;
gabarit FIBAA « Description of course unit according to the ECTS Users' Guide 2015 »
(\url{https://www.fibaa.org/fileadmin/redakteur/pdf/PROG/Handreichungen_und_Vorlagen/DESCRIPTION_Course_Unit_ECTS_User_s_Guide_2015_202004.pdf})~;
UCLouvain, fiche descriptive d'unité d'enseignement
(\url{https://uclouvain.be/fr/system/files/uclouvain_assetmanager/groups/cms-editors-lll/Fiche_Descriptive_0.pdf})~;
ULB, manuel d'édition des fiches UE
(\url{https://ulb.be/fr/programme/manuel-dedition-des-fiches-ue}).

\section{Cible fonctionnelle}\label{sec:cible}

\subsection{Types d'EP définis par la DEVE, sans programmation}\label{sec:types}

Les six types codés en dur deviennent des \textbf{types paramétrés}~: la DEVE en ajoute,
en renomme, en désactive, en supprime (s'il ne porte aucun crédit) et en ordonne
librement. Le comportement d'un type n'est plus déduit de son code mais de ses
paramètres~:
\begin{itemize}
  \item \textbf{Mode d'acquisition}~: \emph{catalogue} (inscription à un EP du catalogue,
    qui porte ses ECTS), \emph{déclaration} (l'étudiant déclare, avec justificatifs),
    \emph{automatique depuis les stages} (barème d'ECTS par jour, lu dans \code{mod_stage}),
    \emph{saisie DEVE} (seule la DEVE crée des crédits de ce type, à la main ou par import).
  \item \textbf{Circuit de validation}, suite ordonnée d'étapes parmi~: \emph{aucune}
    (validé d'office à l'inscription ou à la déclaration), \emph{responsable de l'EP},
    \emph{enseignant référent}, \emph{DEVE}. Exemples~: « responsable » (EP interne
    actuel), « référent » (déclarations actuelles), « référent puis DEVE », « DEVE seule »,
    « aucune ».
  \item \textbf{Évaluation}~: \emph{sans note} (validé / non validé), \emph{avec note}
    (barème et seuil de validation, par défaut sur 20 avec seuil 10), \emph{validé
    d'office} (toute inscription confirmée est validée). Un EP du catalogue peut préciser son
    propre barème dans les limites du type.
  \item \textbf{Justificatifs}~: facultatifs, obligatoires (au moins un fichier), ou sans
    objet~; formats admis (PDF seul, ou documents et images).
  \item \textbf{ECTS}~: fixés par l'EP du catalogue, proposés par l'étudiant (avec un
    maximum), ou calculés (barème par jour)~; plafonds par année et par cursus, inchangés.
  \item \textbf{Supplément au diplôme}~: le type figure ou non en rubrique~4.3.
  \item \textbf{Consigne} affichée à l'étudiant, \textbf{libellé} et \textbf{code court}
    (utilisé dans les imports et les exports).
\end{itemize}
Les six types existants sont migrés automatiquement vers ces paramètres à la mise à jour
(\code{db/upgrade.php}), à l'identique de leur comportement actuel~; les constantes
\code{EP_TYPE_*} ne servent plus qu'à cette migration. L'attribution depuis les stages est
portée par le mode d'acquisition, non plus par le code \code{stage}, ce qui permet
plusieurs barèmes (par exemple un type pour les stages à l'étranger).

\subsection{Catalogue, inscriptions, notes}\label{sec:catalogue}

Un EP du catalogue reçoit~: un \textbf{code}, une \textbf{période d'inscription}
(ouverture et fermeture, facultatives), des \textbf{dates de déroulement}, un
\textbf{état} (brouillon, proposé, publié, clos, archivé), son \textbf{syllabus}
(§\ref{sec:syllabus}) et ses \textbf{prérequis} (§\ref{sec:prerequis}).

Une page \textbf{Participants} par EP, ouverte au responsable et à la DEVE, permet~:
\begin{itemize}
  \item d'\textbf{ajouter} un étudiant après la période d'inscription (au-delà des places
    si la DEVE le décide, avec avertissement), un par un ou par sélection multiple parmi les
    inscrits au cours~;
  \item de \textbf{retirer} un étudiant (motif conservé, place libérée, état « retiré »
    distinct de « refusé »)~;
  \item de \textbf{saisir les notes} ou les validations de toute la liste en une fois
    (tableau éditable, enregistrement en une transaction sous le verrou de l'activité)~; la
    note décide du statut selon le seuil, un « absent » ou « non validé » explicite est
    possible~; les ECTS retenus suivent la validation~;
  \item d'\textbf{importer les notes} par CSV (adresse ou identifiant, note ou
    validé/non validé, commentaire) et d'\textbf{exporter la liste} (CSV, feuille
    d'émargement PDF)~;
  \item de \textbf{clore} l'EP~: les inscriptions restées en attente sont tranchées
    (validées d'office, non validées, ou laissées) selon le choix du responsable.
\end{itemize}
L'inscription de l'étudiant vérifie les places \textbf{sous le verrou de l'activité} et
relit le compte avant d'insérer (défaut~1 de §\ref{sec:defauts}). Les décisions ne sont
écrites que si le crédit est encore dans l'état lu (défaut~3), comme
\code{stage_update_entry_if()} dans StageCompagnon.

\subsection{Proposition d'EP par un enseignant et syllabus}\label{sec:syllabus}

Nouvelle capacité \code{mod/ep:propose} (archétypes enseignant et enseignant éditeur).
Un enseignant peut~:
\begin{itemize}
  \item \textbf{proposer un nouvel EP}~: il crée une fiche à l'état « brouillon », la
    complète, la \textbf{soumet}~; la DEVE la \textbf{publie} (elle entre au catalogue),
    la \textbf{renvoie} avec commentaire, ou la \textbf{refuse}~; chaque passage est
    journalisé et notifié par courriel~;
  \item \textbf{mettre à jour son EP} déjà publié~: les modifications du syllabus sont
    enregistrées comme \textbf{révision en attente}, publiée par la DEVE (ou d'office si
    la DEVE a choisi « publication directe par le responsable » dans les paramètres). Le
    catalogue montre la version publiée tant que la révision n'est pas acceptée.
\end{itemize}
Le \textbf{syllabus} comprend les rubriques ECTS~: objectifs, acquis d'apprentissage,
prérequis (texte libre, plus les prérequis structurés de §\ref{sec:prerequis}), contenu,
méthodes d'enseignement, langue, charge de travail (heures en présence et de travail
personnel, par nature~: cours, TD, TP, terrain, projet), \textbf{conditions d'examen}
(modalités, critères, seuil, seconde session), bibliographie, et la \textbf{liste des
séances}~: pour chacune, intitulé, nature (CM, TD, TP, clinique, autre), durée, intervenant
(compte Moodle ou nom libre pour un intervenant extérieur), date facultative. Un même
enseignant peut intervenir sur plusieurs séances~; le volume horaire total et la
répartition par nature sont calculés.

\textbf{Éditions PDF}~: la fiche complète d'un EP (syllabus, séances, responsables,
modalités), et le \textbf{livret des EP} de la promotion (toutes les fiches publiées, avec
sommaire), générés avec la bibliothèque PDF de Moodle (TCPDF, déjà employée pour le bilan
de promotion de StageCompagnon). Un gabarit de page (logo, en-tête, pied de page) est
réglé par la DEVE.

\subsection{Supplément au diplôme}\label{sec:supplement}

\paragraph{Paramétrage (DEVE, panneau d'administration).} Les textes fixes des rubriques
2 à 5 et 8, le logo, le nom et la qualité du signataire, l'image du cachet, la langue
d'édition (français, anglais facultatif champ par champ), la liste des \textbf{catégories}
d'éléments que l'étudiant peut proposer (par exemple engagement associatif, mobilité,
distinction, certification, publication), et pour ces éléments le \textbf{circuit de
validation} (aucune, enseignant référent, DEVE, référent puis DEVE) et l'\textbf{exigence
de preuve} (PDF obligatoire ou non, nombre maximal, taille). Ce paramétrage est
transférable (§\ref{sec:transfert}).

\paragraph{Contenu automatique.} Tous les crédits validés dont le type est marqué « figure
au supplément » sont portés en rubrique~4.3~: intitulé, année d'étude, ECTS retenus, note
ou « validé », date. L'étudiant ne peut pas en retirer.

\paragraph{Éléments demandés par l'étudiant.} Depuis sa page, l'étudiant \textbf{propose
un élément}~: catégorie, intitulé, description, période, preuves (PDF, enregistrés dans une
zone de fichiers dédiée, conservés avec le dossier). L'élément suit le circuit configuré~;
le validateur voit les preuves, retient ou refuse avec motif~; l'étudiant est notifié.
L'élément validé est porté en rubrique~6.1. L'étudiant peut retirer une demande tant
qu'elle n'est pas tranchée.

\paragraph{Émission.} La DEVE \textbf{émet} le supplément d'un étudiant~: le PDF est
généré, numéroté (numéro unique par activité et année), horodaté, signé du nom et de la
qualité du signataire, porte le cachet, un \textbf{code de vérification} (empreinte
SHA-256 du fichier, affichée sur le document) et la mention « Document établi
conformément au modèle européen ». Le fichier est \textbf{conservé} (zone de fichiers,
ligne d'émission avec auteur, date, empreinte)~: une réédition produit un nouveau numéro et
l'ancien est marqué remplacé. L'étudiant télécharge ses suppléments émis~; avant émission,
il voit un \textbf{aperçu} filigrané « non certifié ». L'\textbf{émission en lot} pour
toute la promotion produit une archive ZIP (tâche ad hoc au-delà de quelques dizaines
d'étudiants, comme le téléchargement groupé des rapports de StageCompagnon), avec un
compte rendu des étudiants dont il manque des données (date de naissance, identifiant).

\paragraph{Données de la rubrique 1.} Nom et prénom du profil Moodle~; identifiant
\code{idnumber}~; date de naissance lue dans un champ de profil personnalisé désigné par la
DEVE, à défaut dans le détail de convention de \code{mod_stage} (\code{birthdate}), à
défaut laissée à compléter avant émission (l'émission d'un supplément incomplet est
refusée, sauf case « émettre malgré les champs manquants »).

\subsection{Prérequis d'EP et validation de l'année}\label{sec:prerequis}

Un EP du catalogue peut exiger, pour l'inscription~: un ou plusieurs \textbf{EP validés}
(prérequis structurés), une \textbf{année d'étude minimale} (existe), un \textbf{nombre
minimal d'ECTS} déjà retenus. Le catalogue affiche à l'étudiant ce qui lui manque~; la
règle est bloquante ou seulement signalée, au choix de la DEVE par EP~; la DEVE peut
inscrire malgré tout depuis la page Participants.

La \textbf{validation de l'année} devient une procédure explicite~: pour chaque étudiant
et chaque année d'étude, la DEVE constate (en lot, depuis le pilotage, avec filtres « à
jour / en défaut ») que les minimums sont atteints, et enregistre « année validée » (date,
auteur, commentaire), ou « année non validée » avec motif. Le bilan de l'étudiant, les
exports et le supplément s'appuient sur cette décision. La \textbf{rentrée} est une
procédure guidée~: passage de l'année courante à la suivante, rappel des étudiants en
défaut, et \textbf{transfert d'un étudiant} (redoublement, changement de promotion) vers
une autre activité \code{mod_ep} du site, crédits, preuves et éléments de supplément
compris, sur le modèle du transfert de StageCompagnon.

\subsection{Imports, exports, transfert du paramétrage}\label{sec:transfert}

\paragraph{Import CSV des EP (DEVE).} Deux fichiers, avec modèle téléchargeable,
documentation des colonnes dans l'interface et dans \code{INSTALL.md}, prévisualisation
ligne par ligne (erreurs, doublons, lignes ignorées), confirmation sous verrou et en une
transaction~:
\begin{itemize}
  \item \textbf{catalogue}~: code, intitulé, type (code), ECTS, années, places, période,
    responsables (adresses), rubriques du syllabus, séances (une ligne par séance avec le
    code de l'EP)~;
  \item \textbf{crédits}~: étudiant (adresse ou identifiant), type (code), EP du catalogue
    (code, facultatif), intitulé, année d'étude, ECTS demandés et retenus, note, statut,
    date, validateur (adresse), commentaire. Sert aussi à reprendre l'historique d'une
    promotion.
\end{itemize}

\paragraph{Export global du module.} Un classeur Excel avec une feuille par table
fonctionnelle~: paramètres, types, catalogue et syllabus, séances, responsables,
prérequis, minimums, bilan par étudiant, crédits, éléments de supplément, suppléments émis.
Toutes les colonnes nécessaires à une \textbf{restauration} y figurent, et
\code{import_global.php} relit ce classeur dans une activité vide (ou en complément), avec
la même prévisualisation que les imports CSV. Les fichiers (justificatifs, preuves, PDF
émis) ne tiennent pas dans un classeur~: une archive ZIP les accompagne, avec un index
reliant chaque fichier à sa ligne.

\paragraph{Transfert du paramétrage.} Deux voies, l'une n'excluant pas l'autre~:
\begin{itemize}
  \item \textbf{export et import du paramétrage} en un fichier JSON signé par version
    (types, minimums, catégories et textes du supplément, gabarit PDF, réglages de
    proposition), depuis le panneau d'administration~;
  \item \textbf{copie depuis une autre activité} du site, choisie dans une liste (les
    activités où l'utilisateur a \code{mod/ep:manage}), avec cases pour choisir ce qui est
    copié, y compris le catalogue et ses syllabus (sans inscriptions).
\end{itemize}
La sauvegarde Moodle (§\ref{sec:backup}) emporte en outre tout cela lors d'une
duplication ou d'un import de cours.

\subsection{Sauvegarde, journaux, données personnelles}\label{sec:backup}

\begin{itemize}
  \item \code{backup/moodle2/} pour \code{mod_ep}~: paramétrage toujours sauvegardé (types,
    catalogue, syllabus, séances, responsables, prérequis, minimums, textes et images du
    supplément, gabarit PDF)~; avec les données utilisateur~: crédits, notes, validations
    d'année, éléments de supplément, émissions, et tous les fichiers. Pour
    \code{mod_epsynthesis}~: liens, conservés sur le même site. Tests de sauvegarde et
    restauration comme dans StageCompagnon. Dates non décalées à la restauration.
  \item \textbf{Événements Moodle} pour chaque transition~: inscription, ajout et retrait
    par le responsable, note saisie, validation, refus, annulation, déclaration créée et
    modifiée, proposition soumise, renvoyée, publiée, refusée, révision publiée, élément de
    supplément demandé, validé, refusé, supplément émis, année validée, transfert, import
    (avec le nombre de lignes).
  \item \textbf{Fournisseur RGPD} étendu aux nouvelles tables et zones de fichiers~; même
    règle que StageCompagnon~: l'étudiant supprimé emporte tout son dossier (preuves et
    suppléments émis compris, avec l'avertissement d'archivage dans \code{INSTALL.md}), le
    personnel supprimé est seulement dissocié.
  \item \textbf{Notifications} par courriel, avec modèles modifiables par la DEVE (comme les
    courriels de StageCompagnon)~: décision sur une demande, proposition à examiner,
    supplément émis, relance des demandes en attente (tâche planifiée hebdomadaire).
  \item \textbf{Réinitialisation de cours}~: options pour supprimer les crédits, les
    éléments de supplément et les émissions de l'année écoulée.
\end{itemize}

\subsection{Synthèse enseignant}

\code{mod_epsynthesis} suit~: les propositions et révisions de l'enseignant et leur état,
les EP dont il est responsable avec le nombre d'inscrits à noter, les éléments de
supplément de ses étudiants à valider, en plus des files existantes. Il reçoit sa
sauvegarde Moodle, ses événements et ses tests.

\section{Modèle de données cible}

Les tables actuelles sont conservées et étendues~; les nouvelles sont ajoutées par
\code{db/upgrade.php} avec points de sauvegarde (contrôle \code{savepoints} de la CI).

\begin{longtable}{p{4.6cm}p{10.3cm}}
\toprule
\textbf{Table} & \textbf{Évolution} \\
\midrule
\endhead
\code{ep} & \code{proposalsdirectpublish}, \code{birthdatefield},
  \code{supplementflow}, \code{supplementevidencerequired}, \code{supplementnumbering},
  \code{pdfheader}, \code{pdffooter}, \code{timesyncedlock} \\
\code{ep_type} & \code{acquisition}, \code{validationflow}, \code{grading},
  \code{grademax}, \code{gradepass}, \code{evidence}, \code{evidencetypes},
  \code{ectsmode}, \code{maxclaimedects}, \code{insupplement}~; code libre de 2 à 20
  caractères \\
\code{ep_activity} & \code{code}, \code{state}, \code{registrationstart},
  \code{registrationend}, \code{datestart}, \code{dateend}, \code{grading},
  \code{grademax}, \code{gradepass}, \code{closeaction}, \code{proposedby},
  \code{publishedby}, \code{timepublished}, \code{prerequisitemode} \\
\code{ep_syllabus} (nouvelle) & une ligne par EP et par version (publiée, en attente)~:
  rubriques ECTS en texte formaté, charge de travail par nature, langue, bibliographie,
  conditions d'examen, \code{status}, \code{reviewcomment}, auteur, dates \\
\code{ep_session} (nouvelle) & séances d'un syllabus~: intitulé, nature, durée, date,
  \code{teacherid} ou \code{teachername}, ordre \\
\code{ep_prerequisite} (nouvelle) & EP requis, année minimale, ECTS minimaux,
  \code{blocking} \\
\code{ep_credit} & \code{grade}, \code{gradedby}, \code{timegraded}, \code{step} (étape
  courante du circuit), \code{addedby} (ajout par le responsable), index unique
  \code{epid, source, sourceref} pour les crédits automatiques, index
  \code{epid, status}, \code{activityid, status}, \code{userid} \\
\code{ep_credit_decision} (nouvelle) & historique des décisions par étape (qui, quand,
  quoi, commentaire), pour les circuits à plusieurs étapes et la traçabilité \\
\code{ep_year_validation} (nouvelle) & étudiant, année, décision, auteur, date,
  commentaire \\
\code{ep_supplement_category} (nouvelle) & catégories d'éléments proposables, ordre,
  actif \\
\code{ep_supplement_item} (nouvelle) & élément demandé par l'étudiant, statut, étape,
  validateur, commentaire, zone de fichiers \code{supplementevidence} \\
\code{ep_supplement_issue} (nouvelle) & émission~: numéro, étudiant, auteur, date,
  empreinte, langue, \code{replacedby}, fichier en zone \code{supplementissued} \\
\code{ep_supplement_text} (nouvelle) & textes des rubriques fixes, par langue \\
\code{ep_notification} (nouvelle) & modèles de courriels modifiables \\
\code{ep_import_log} (nouvelle) & imports réalisés (fichier, auteur, lignes, date), pour
  retrouver et défaire un import \\
\bottomrule
\end{longtable}

\section{Efficacité et accès simultanés}\label{sec:perf}

Objectif chiffré, mesuré sur le banc (§\ref{sec:banc})~: toute page de liste en moins
d'une seconde et moins de 40~requêtes pour une promotion de 250 étudiants et 100\,000
crédits sur le site~; export global de 20 promotions en moins de deux minutes.
\begin{itemize}
  \item \textbf{Pagination en SQL} des listes (validation, pilotage, participants,
    catalogue) avec comptage séparé~; recherche et tri en base.
  \item \textbf{Cache de bilan de promotion} (classe \code{progress_cache}, comme
    StageCompagnon)~: types, minimums, crédits de tous les étudiants de la page chargés en
    trois requêtes~; \code{ep_get_student_progress()} s'y alimente quand il est armé.
  \item \textbf{Chargements groupés}~: responsables, places prises, inscriptions de
    l'étudiant courant, fichiers de preuve, en une requête par page et non par ligne.
  \item \textbf{Index} listés dans le modèle de données, vérifiés par \code{EXPLAIN
    ANALYZE} sur la base de stress.
  \item \textbf{Verrou par activité} (\code{ep_get_activity_lock()}, fabrique de verrous
    Moodle) sur~: inscription, ajout de participants, saisie des notes en lot, imports,
    synchronisation des stages, émission en lot, transfert.
  \item \textbf{Écritures conditionnées} à l'état lu (\code{ep_update_credit_if()}) pour
    toute décision~; message « cette demande a changé entre-temps » à l'utilisateur dont
    l'écriture n'a pas eu lieu.
  \item \textbf{Synchronisation des stages}~: sous verrou, et mise à jour de
    \code{timesynced} par requête conditionnelle avant le calcul, pour qu'une seule page
    déclenche la synchronisation.
  \item \textbf{Tâches ad hoc} pour les traitements longs (émission en lot, ZIP d'export,
    notifications), avec compte rendu par courriel.
  \item \textbf{Gabarits Mustache} et classes \code{output} pour les listes, badges,
    filtres et actions, partagés avec \code{mod_epsynthesis}.
\end{itemize}

\section{Sécurité}

Revue systématique avant livraison de chaque lot, et rapport dans \code{docs/audit/}~:
\begin{itemize}
  \item chaque action d'écriture passe par \code{require_sesskey()} ou
    \code{confirm_sesskey()} et une capacité vérifiée dans le contexte du module~; les
    identifiants reçus sont toujours contraints à l'instance (\code{epid}) et au périmètre
    de l'utilisateur (responsable, référent, propriétaire)~;
  \item accès aux fichiers (preuves, PDF émis) par un script de contrôle, comme
    \code{evidence_file.php}, jamais par \code{pluginfile} ouvert~;
  \item \code{PARAM_*} exacts (fin des \code{PARAM_RAW})~; \code{returnurl} en
    \code{PARAM_LOCALURL}~; texte formaté passé par \code{format_text} avec le format
    enregistré~; plus de JavaScript en ligne construit à partir de chaînes~;
  \item imports~: taille et type de fichier contrôlés, cellules nettoyées, formules Excel
    neutralisées à l'export (préfixe d'apostrophe, comme StageCompagnon)~;
  \item PDF émis~: empreinte vérifiable, fichier en lecture seule, numéro non devinable
    (année, séquence et contrôle)~;
  \item tests PHPUnit « frontières de droits » pour chaque page (un étudiant sur le crédit
    d'un autre, un responsable sur un EP qui n'est pas le sien, un référent sur un étudiant
    qui n'est pas le sien, une activité sur les données d'une autre)~;
  \item passage de la revue de sécurité outillée (\code{/security-review}) sur chaque lot et
    de l'analyse statique Moodle de la CI.
\end{itemize}

\section{Tests et contrôle continu}\label{sec:tests}

\subsection{Tests unitaires et de scénario}

\begin{itemize}
  \item \textbf{PHPUnit}~: un fichier par règle métier (types paramétrés, circuits de
    validation, notes et seuils, places sous concurrence simulée, prérequis, validation
    d'année, supplément et émission, imports, export et restauration, sauvegarde et
    restauration Moodle, RGPD, événements, transfert)~; données tabulaires en
    \code{dataProvider}~; les tests existants sont conservés et adaptés.
  \item \textbf{Behat}, sans JavaScript, sur le serveur intégré de PHP, un scénario par
    procédure~: la DEVE paramètre un type et le catalogue, un enseignant propose un EP et la
    DEVE le publie, un étudiant s'inscrit, le responsable ajoute un étudiant et saisit les
    notes, un étudiant déclare un EP et le référent le valide, un étudiant demande un
    élément de supplément avec preuve PDF et la DEVE l'accepte puis émet le supplément,
    la DEVE importe un CSV, exporte et restaure le module, valide l'année.
  \item \textbf{CI GitHub} reprise de StageCompagnon~: \code{phpcs} standard Moodle (même
    \code{phpcs.xml}), \code{moodle-plugin-ci} sur Moodle 4.3 / PHP 8.0 / PostgreSQL,
    4.5 / 8.2 / PostgreSQL avec couverture et Behat, 4.5 / 8.2 / MariaDB~; lint,
    validation, savepoints, PHPDoc, Mustache bloquants. \code{mod_stage} est installé en
    plugin supplémentaire puisque \code{mod_ep} en dépend.
\end{itemize}

\subsection{Banc dans le cloud}\label{sec:banc}

L'environnement de cette session dispose de PHP~8.3, PostgreSQL~16 et Docker~: un Moodle
4.5 réel y est installé avec \code{mod_stage}, \code{mod_ep} et leurs synthèses, comme
l'a fait l'audit de concurrence de StageCompagnon. Scripts versionnés dans
\code{docs/audit/banc/}~:
\begin{itemize}
  \item \textbf{Volumétrie}~: 20 promotions, 5\,000 étudiants, 400 enseignants, 300 EP au
    catalogue avec syllabus, 100\,000 crédits, 20\,000 éléments de supplément, 50\,000
    fichiers de preuve~; stages complémentaires dans \code{mod_stage} pour l'attribution
    automatique.
  \item \textbf{Chronométrage de toutes les pages} pour chaque rôle, nombre de requêtes et
    mémoire, avant et après correction, avec \code{EXPLAIN ANALYZE} des requêtes lentes.
  \item \textbf{Concurrence}~: 50 étudiants sur les 5 dernières places d'un EP~; un
    responsable et la DEVE statuant sur les mêmes crédits~; deux DEVE confirmant le même
    import~; synchronisation des stages pendant une validation DEVE dans \code{mod_stage}~;
    émission en lot pendant qu'un étudiant ajoute un élément~; vérification des invariants
    (pas de surbooking, pas de crédit en double, pas de décision perdue, pas de verrou
    restant).
  \item \textbf{Procédures de bout en bout} jouées par scripts HTTP (sessions par rôle)~:
    DEVE, enseignant référent, responsable d'EP, proposant, étudiant, jusqu'à la
    validation de l'année et l'émission du supplément de la promotion.
  \item Rapports LaTeX \code{docs/audit/test-de-charge.tex} et
    \code{docs/audit/test-de-concurrence.tex}, avec mesures avant et après.
\end{itemize}

\section{Documentation}

\code{README.md} et \code{INSTALL.md} des deux modules au niveau de StageCompagnon
(rôles, mise en place, chaque procédure, colonnes des imports, RGPD, sauvegarde)~;
\code{CHANGELOG.md} tenu à chaque changement de \code{version.php}~; trois guides LaTeX
avec captures d'écran prises par script (\code{docs/guides/})~: paramétrage DEVE (types,
catalogue, supplément, imports, rentrée), enseignant (référent, responsable d'EP,
proposant), étudiant (inscription, déclaration, supplément). Les PDF sont compilés et
versionnés avec leurs sources, comme dans StageCompagnon.

\section{Lots et ordre de réalisation}\label{sec:lots}

Chaque lot se termine par ses tests PHPUnit, le passage de la CI, une entrée de
\code{CHANGELOG.md}, un numéro de version et une poussée sur la branche. L'ordre suit les
dépendances~: les fondations (lot~1) servent à tout le reste, le modèle de types (lot~2)
conditionne les circuits (lots~3 à 5), les imports et exports (lot~6) ont besoin de toutes
les tables, le banc (lot~8) se joue sur le produit complet.

\begin{longtable}{p{1.1cm}p{5.2cm}p{5.2cm}p{3.2cm}}
\toprule
\textbf{Lot} & \textbf{Contenu} & \textbf{Livrables} & \textbf{Dépend de} \\
\midrule
\endhead
1 & Fondations~: CI et \code{phpcs}, Moodle 4.3+, verrou d'activité, écritures
  conditionnées, index, sauvegarde \code{backup/moodle2}, événements, Mustache, correction
  des défauts de §\ref{sec:defauts}, cache de bilan, pagination SQL &
  \code{.github/workflows/ci.yml}, \code{phpcs.xml}, \code{backup/}, \code{classes/event/},
  \code{templates/}, tests de sauvegarde, d'événements et de concurrence simulée & -- \\
2 & Types paramétrés~: modèle, migration des six types, page d'administration (ajout,
  édition, ordre, suppression), circuits de validation à étapes, historique des décisions,
  justificatifs obligatoires & \code{types.php} refondu, \code{db/upgrade.php},
  \code{ep_credit_decision}, tests de circuits & 1 \\
3 & Catalogue et participants~: états et périodes de l'EP, page Participants (ajout,
  retrait, notes en lot, import des notes, clôture), inscription sous verrou, prérequis
  structurés & \code{participants.php}, \code{grades_import.php}, tests de notes, de
  places et de prérequis, Behat « inscription et notes » & 2 \\
4 & Proposition d'EP et syllabus~: fiche, séances, intervenants, soumission, revue,
  publication, révisions~; PDF de fiche et livret~; notifications & \code{propose.php},
  \code{syllabus.php}, \code{classes/pdf/}, modèles de courriels, Behat « proposition » & 2 \\
5 & Supplément au diplôme~: paramétrage, catégories, éléments de l'étudiant avec preuves
  PDF, circuit, aperçu, émission unitaire et en lot (ZIP, tâche ad hoc), vérification,
  validation d'année et rentrée, transfert d'étudiant & \code{supplement*.php},
  \code{year_validation.php}, \code{transfer.php}, tests, Behat « supplément » & 2, 3 \\
6 & Imports et exports~: CSV catalogue et crédits avec prévisualisation, export global
  Excel et ZIP, restauration d'export global, export et import JSON du paramétrage, copie
  depuis une autre activité, journal des imports & \code{import*.php},
  \code{export_global.php}, \code{classes/local/*importer.php}, tests & 2 à 5 \\
7 & Synthèse enseignant et RGPD~: \code{mod_epsynthesis} étendu, sauvegarde, événements~;
  fournisseur RGPD des deux modules, réinitialisation de cours & tests RGPD et de synthèse
  & 4, 5 \\
8 & Banc cloud~: installation Moodle réel, jeux de données, chronométrage, concurrence,
  procédures de bout en bout, corrections, rapports LaTeX & \code{docs/audit/} & 1 à 7 \\
9 & Documentation et sécurité finale~: README, INSTALL, CHANGELOG, guides avec captures,
  revue de sécurité, version \code{1.0.0} & \code{docs/guides/}, rapport de sécurité & 1 à 8 \\
\bottomrule
\end{longtable}

Le lot~1 est engagé en premier parce qu'il corrige des défauts qui existent aujourd'hui
et qu'il fournit l'outillage (CI, verrou, événements, sauvegarde) sur lequel chaque lot
suivant s'appuie pour être testé.

\section{Hypothèses retenues et questions}\label{sec:questions}

Pour ne pas bloquer le démarrage, les choix suivants sont faits par défaut~; chacun
peut être inversé sans remettre en cause le plan.
\begin{enumerate}
  \item \textbf{Barème par défaut}~: note sur 20, seuil 10, pas de demi-point imposé~;
    modifiable par type et par EP. Les notes ne sont \textbf{pas} poussées dans le carnet
    de notes Moodle (\code{FEATURE_GRADE_HAS_GRADE} reste à faux)~: le carnet d'un cours de
    promotion n'est pas utilisé pour les EP. À confirmer.
  \item \textbf{Qui propose un EP}~: tout utilisateur ayant \code{mod/ep:propose} dans le
    cours de la promotion (archétypes enseignant et enseignant éditeur). Un enseignant
    extérieur au cours doit y être inscrit (rôle enseignant non éditeur) pour proposer.
  \item \textbf{Date de naissance} pour la rubrique~1.3~: champ de profil personnalisé
    choisi par la DEVE, puis \code{mod_stage}, puis saisie manuelle à l'émission.
  \item \textbf{Langue du supplément}~: français~; chaque texte fixe peut recevoir une
    version anglaise, et l'émission se fait dans l'une ou l'autre langue. Pas d'édition
    bilingue en colonnes.
  \item \textbf{Ce qui figure en rubrique 4.3}~: tous les crédits validés des types marqués
    « figure au supplément », sans possibilité pour l'étudiant d'en retirer~; les EP non
    validés et annulés n'y figurent jamais.
  \item \textbf{Certification}~: signature par nom, qualité et cachet image, plus
    empreinte SHA-256 imprimée et conservée. Pas de signature électronique qualifiée
    (certificat) dans cette version~; la structure permet d'en ajouter une plus tard.
  \item \textbf{Séances d'un EP}~: la fiche décrit le programme des séances (nature,
    durée, intervenant)~; elle ne gère ni planning horaire ni salles.
  \item \textbf{Synthèse}~: la décision continue de se prendre dans le cours de la
    promotion, la synthèse n'étant qu'une vue, comme aujourd'hui.
  \item \textbf{Moodle minimal}~: 4.3 et PHP 8.0, alignés sur StageCompagnon et sur sa CI.
\end{enumerate}

\end{document}
